\subsection{Введение в проблематику}

Современное машиностроительное производство невозможно представить без использования систем автоматизированного проектирования (САПР). Такие системы, как КОМПАС-3D, SolidWorks, CATIA, NX и другие, стали основным инструментом конструктора, позволяя создавать сложные трёхмерные модели деталей и сборочных единиц с высокой точностью и производительностью. Ключевым преимуществом современных САПР является параметрическое моделирование — подход, при котором геометрия модели описывается не фиксированным набором точек и поверхностей, а последовательностью операций с изменяемыми числовыми параметрами и геометрическими ограничениями.

Однако на практике инженеры регулярно сталкиваются с ситуацией, когда параметрическая история построения модели недоступна. Это происходит при импорте моделей из нейтральных форматов (например, STEP или IGES), при получении моделей от сторонних организаций, использующих другое программное обеспечение, при работе с унаследованными базами моделей, созданных без соблюдения параметрической дисциплины, а также при получении геометрии с 3D-сканеров реальных физических объектов \cite{kim2023reconstruction}. В таких случаях инженер получает «немую» геометрию — статичное представление, из которого невозможно понять логику создания модели, невозможно изменить отдельные элементы, не перестраивая модель заново. Редактирование такой модели становится крайне трудоёмким, а часто и вовсе невозможным без полного перепроектирования.

Задача восстановления истории построения по финальной геометрии, известная как обратный инжиниринг конструктивной истории (Inverse CAD), является одной из наиболее сложных и востребованных задач в области промышленного программного обеспечения. Её решение позволило бы не только понимать логику создания изделия, но и эффективно адаптировать его под новые требования, исправлять ошибки, использовать существующие разработки в качестве основы для новых проектов, а также обучать начинающих специалистов на примерах реальных моделей \cite{autodesk2020}.

Традиционные алгоритмические методы анализа геометрии, основанные на поиске примитивов и распознавании поверхностей, имеют существенные ограничения. Одна и та же итоговая форма может быть получена различными последовательностями операций, а одно и то же локальное изменение геометрии может иметь множество параметрических интерпретаций \cite{isicad2013direct}. Более того, такие методы, как правило, производят искусственный, не свойственный человеку порядок операций, что затрудняет их практическое использование.

В последние годы методы машинного обучения, особенно глубокие нейронные сети, показали впечатляющие результаты в анализе трёхмерных данных (классификация, сегментация, генерация). Однако прямое применение существующих подходов для восстановления именно истории построения наталкивается на ряд фундаментальных ограничений, связанных с необходимостью предсказывать не только тип операции, но и её точные числовые параметры, чувствительные к масштабу, системе координат и дискретизации.

Данная глава посвящена систематическому анализу теоретических и технологических основ, необходимых для решения задачи восстановления истории построения CAD-моделей с использованием методов машинного обучения. В ней будут рассмотрены: принципы параметрического моделирования и структура истории построения, форматы представления трёхмерных данных (STEP и STL), современное состояние методов обратного инжиниринга в промышленности, а также обзор архитектур нейронных сетей для анализа трёхмерной геометрии — от облаков точек до сеточных представлений.

